\section{Metodologia Experimental}
\label{sec:setup}

\subsection{Bases}
\textbf{ASEADOS-SDN-IoT}~\cite{aseados}: registros de fluxo de uma rede SDN-IoT. Usamos as classes
Benign e DDoS (312.262 fluxos, 16,5\% de ataque), como o artigo, e todas as cinco classes (Benign,
Bot, DDoS, DoS, Probe; 457.044 fluxos) para o teste com ataques não-DDoS. A ordem é a do artigo
(Seção~\ref{sec:replication}), em que o DDoS dura 1.117\,s, a cerca de 46 fluxos/s, depois de
233 mil fluxos benignos.
\textbf{CTU-13, cenário 11}~\cite{garcia2014ctu13}: 6.336.394 pacotes em 971\,s (cerca de
6.500 pacotes/s), com o \emph{botnet} Rbot fazendo um \emph{flood} ICMP de cerca de 17 mil
pacotes/s contra um alvo. O rótulo é a regra oficial: o pacote é ataque se envolve um dos três IPs
infectados (32,0\%).
\textbf{Edge-IIoTset}~\cite{ferrag2022edgeiiot}: \emph{pcaps} de um \emph{testbed} IIoT, todos os
quadros, ordenados por tempo.

\subsection{Configurações}
Para não ajustar hiperparâmetros na base de teste, todas as rodadas do \modusp\ usam uma única
configuração, \texttt{fixed\_q}, definida na Etapa~1:
HST com os valores padrão de Tan \emph{et al.}~\cite{tan2011hst} ($W=250$, 25 árvores, profundidade
15, \emph{size limit} 0,1); limiar $\tilde{S} + 5{,}19\cdot\mathrm{MAD}$ (z-score modificado de
Iglewicz e Hoaglin) sobre $|S|=5.000$; piso automático igual ao quantil 0,999 dos scores do
aquecimento; e ADWIN com $\delta=0{,}001$ no tamanho. Os parâmetros novos foram fixados antes das
rodadas: janela de 10\,s, teto por destino no percentil 99 do aquecimento, quantil 0,999 para o
limiar congelado, $B_{\min}=B_{\mathrm{est}}=5.000$ e $B_{\max}=50.000$. As variações desses
valores aparecem apenas como análise de sensibilidade (Seção~\ref{sec:sens}), sempre identificadas
como tal. Como referência, reportamos também o \modus\ com os hiperparâmetros do artigo para cada
base.

\subsection{Métricas}
Reportamos precisão, recall e F1 da classe ataque na fase operacional, isto é, após o fim do
aquecimento de cada rodada. Para comparar aquecimentos de tamanhos diferentes, também calculamos as
métricas no trecho comum após as 50 mil instâncias. Reportamos ainda duas métricas independentes
de limiar: a AUC-ROC e a PR-AUC dos scores (no \modusp, do modelo adaptativo; quando indicado, do
congelado). Para medir o envenenamento diretamente, reportamos a fração das instâncias de ataque
pós-aquecimento que foram usadas no treino do modelo adaptativo.

\subsection{Implementação e reprodutibilidade}
O código, os \emph{logs} e as métricas de todas as rodadas estão no repositório do trabalho. Os
\emph{scores} por instância são regenerados pelo \texttt{run.py}. As rodadas foram feitas em um Mac
ARM de 10 núcleos e 16\,GB, com Python 3.13, CapyMOA 0.15 e OpenJDK 17, várias em paralelo. Por isso
os tempos são indicativos.
